\honest{The Sin of Underconfidence}{Eliezer Yudkowsky}{2009}

There are ``three great besetting sins of rationalists,'' and the third is underconfidence. I do not name the other two. Michael Vassar regularly accuses me of it, and is the only person on Earth who does. \nb{``Final Words,'' posted a week later and placed earlier in the book, has Jeffreyssai warn of three flaws, the third of which is underconfidence.}

People warned about a bias sometimes correct too far. You know you are biased but not how much, so you correct a little more, and more, and wonder whether you have overshot, until estimation feels futile. \nb{Research on bias correction agrees: Wegener and Petty write that corrections ``often'' produce a bias in the other direction.} We tend to cast overconfidence as the sin of pride, from that other list that never warned against misused humility, and fear the humiliation of being cast down. Modesty seems lighter: being found better than you thought is a warm surprise, and putting yourself below others seems nice, the sort of thing Gandalf would do. So a reader of many studies of overconfidence may push their self-estimate too low. \nb{In ``Knowing About Biases Can Hurt People'' (2007), Yudkowsky had described another harm from the same knowledge: ammunition against other people's arguments. Here it is turned on oneself.}

The danger is passing up what you could have done. My first advice: if you can find out how good you are, test it. I thought I could win the AI-Box Experiment, and tested it; later, with large sums at stake, I won ``1 time out of 3.'' That was my limit at the time, found without arguing myself up or down.

Smart people who are used to winning stop stretching themselves. ``It is said'' that this comes from defining yourself by intelligence rather than effort. \nb{This is Carol Dweck's research. A close replication in 2019 did not find the classic study's effects on motivation.} I am not sure that is how a rationalist should think, since rationality is systematized winning and trying to try leads to failure. A hypothesis affords testing: if you do not know whether you will win on a hard problem, find out. Congratulating yourself on trying is a bad habit, but not trying is worse. You may tell yourself you have gained information about your level, so long as you add that you will try not to gain the same information next time. Try to win every time, but if you always win, you are not stretching. Console yourself too much and you become a scrub. A master of the Competitive Conspiracy might say: ``It's not okay to lose. But the hurt of losing is not something so scary that you should flee the challenge for fear of it.''

Seek challenges that could humiliate you. I read of a theist who had beaten Christopher Hitchens badly in debate, by atheists' own account, and asked Bloggingheads to arrange a debate with him. Since people said Hitchens should have watched the theist's earlier debates, I decided not to prepare, to learn whether I can ``handle damn near anything on the fly.'' This is not self-handicapping, since I will not treat it as an excuse. \nb{The theist was William Lane Craig. When commenters urged preparation, Yudkowsky wrote the next day that ``it might still be a test of my ability even if I prepare.'' Craig declined to debate.} If I lose, I lose my stake and learn my limits. If you must win, of course, make it as easy as you can; anything else would be spectacular overconfidence, even at tic-tac-toe against a three-year-old.

A subtler form is lost momentum, amid all the things humans do wrong. You become timid; you raise doubts about yourself, never answer or test them, and with no single decision, slow down. Fixing one thing seems pointless when a dozen others stay wrong. Wisdom comes to seem like ever doubting and never resolving, the humility of refusal and never of preparation, and saying worse and worse things about human abilities, down to cynicism.

My last advice is a question to ask of any habit of thought: ``Does this way of thinking make me stronger, or weaker?'' I have warned before against reasonable-sounding arguments, such as two-boxing on Newcomb's problem, doubting all knowledge because of induction, or always adopting the majority belief. \nb{That dispute remains open; the notes on ``Newcomb's Problem and Regret of Rationality'' give the surveys.} Does constantly reminding yourself to doubt make you stronger? Never resolving doubts? A deliberate crisis of faith under uncertainty? Answering every objection with a humble confession of fallibility? If you take more precautions, test yourself more, ask friends, work up to big things gradually and fail less often, you are probably getting stronger. If you never fail, avoid challenges and feel hopeless, you are probably getting weaker.

As a child I found my practice scores falling because I erased correct answers; when I decided to ``use the Force'' and trust my instinct, they rose above where they began, and the real test was higher still. Doubt does not always make you stronger, especially when it blocks good information such as math intuitions, though I needed the test to learn this. \nb{In general, research finds that changed answers are usually changes from wrong to right. The author's case was the exception, found by testing, as the post says.}

Underconfidence is not unique to rationalists, but the attempt to be rational can lead to it. It is a ``stopping mistake'': it prevents the experience that would correct it. Because it ``does seem quite common among aspiring rationalists who I meet,'' though less so among famous role models, I rank it third.
